\documentclass[11pt,a4paper]{article}

\usepackage[T1]{fontenc}
\usepackage[utf8]{inputenc}
\usepackage{lmodern}
\usepackage{microtype}
\usepackage[a4paper,margin=28mm]{geometry}
\usepackage{amsmath,amssymb,mathtools}
\usepackage{booktabs}
\usepackage{array}
\usepackage{longtable}
\usepackage{enumitem}
\usepackage{xcolor}
\usepackage{hyperref}
\usepackage{cleveref}
\usepackage{listings}
\usepackage{url}

\hypersetup{
  colorlinks=true,
  linkcolor=blue!55!black,
  citecolor=blue!55!black,
  urlcolor=blue!60!black,
  pdftitle={Split-Complex Lorentzian Dynamics I: Kinematic Separation, Projective Null Boundary, and the Missing Coupling Structure},
  pdfauthor={Jan Seeck}
}

\newcommand{\R}{\mathbb{R}}
\newcommand{\D}{\mathbb{D}}
\newcommand{\sech}{\operatorname{sech}}
\newcommand{\artanh}{\operatorname{artanh}}
\newcommand{\Stab}{\operatorname{Stab}}
\newcommand{\Orb}{\operatorname{Orb}}
\newcommand{\Rayp}{\operatorname{Ray}_{+}}
\newcommand{\TaskOne}{\textsc{Task 1}}
\newcommand{\TaskTwo}{\textsc{Task 2}}
\newcommand{\newinput}{\textsc{New Input Required}}
\newcommand{\derived}{\textsc{Derived}}
\newcommand{\independent}{\textsc{Independent Degrees of Freedom}}

\title{\textbf{Split-Complex Lorentzian Dynamics I}\\[0.35em]
\large Kinematic Separation, Projective Null Boundary, and the Missing Coupling Structure}
\author{Jan Seeck\\
\small Formal verification: Aristotle (Harmonic)\\
\small AI-assisted analysis and manuscript preparation: ChatGPT (OpenAI)}
\date{31 August 2026}

\begin{document}
\maketitle

\begin{abstract}
This paper opens a new experiment following the completed \emph{Split-Complex Lorentzian Kinematics} dependency audit.  The title ``Dynamics'' is programmatic: no dynamical law is introduced here.  Instead, two machine-checked audits isolate the mathematical structures that would have to be coupled before any dynamics could be claimed.

The internal motivation began from a deliberately non-mathematical visualization involving two counter-propagating periodic carriers and a local periodic process.  That visualization was translated into standard mathematical objects---split-complex null sectors, Lorentz boosts, wave covectors, a timelike worldline, and an independent scalar phase---and was then quarantined from the formal work.  Aristotle received only the mathematical definitions and audit questions; it was not given the heuristic visualization or any intended interpretation.

\TaskOne establishes the standard $1+1$ Lorentz/split-complex structure, derives the wave-covector transformation from phase invariance, proves that a single trajectory-local phase observable identifies only $kv-\omega$, and shows that standard null dispersion fixes the null propagation speed while leaving trajectory-local parameterizations underdetermined before dispersion is imposed.  \TaskTwo refines the null boundary by separating nonzero vectors, positive rays, and projective directions.  It proves that the boost orbit of a scaled null vector is its entire positive ray, whereas quotienting by scale collapses that orbit and enlarges the null stabilizer from trivial to the whole connected boost group.  Independently, a scalar phase $\Theta(\tau)=\Theta_0+\Omega\tau$ satisfies $d\Theta/dt=\Omega\sech\eta$ while its intrinsic proper-parameter rate remains $\Omega$.

The principal result is therefore a missing dependency rather than a new physical law: the Lorentz/null-sector branch and the timelike scalar-phase branch share the rapidity variable but no theorem in the present assumptions relates the positive continuous scale $\Omega$ to the null-wave scale $k_\sigma$.  Any such cross-branch relation is additional structure.  This is the first unresolved node at which a genuine coupling, and hence potentially a dynamics, could enter.
\end{abstract}

\section{Context and scope}

This work follows the completed experiment \href{https://www.astronews.com/community/threads/split-complex-lorentzian-kinematics.12707/post-156281}{\emph{Split-Complex Lorentzian Kinematics}}, in which a split-complex/Lorentz structure was repeatedly decomposed and reconstructed from progressively more elementary mathematical dependencies.  The present experiment asks a different question: once the standard Lorentz/null structure is kept explicit, what additional mathematical structure would be required to relate it to an independent local periodic degree of freedom?

The formalization discussed here was produced in the Aristotle request
\begin{center}
\url{https://aristotle.harmonic.fun/dashboard/requests/201d2dcd-4452-4ae1-8656-4637890c2223}.
\end{center}
The two audited layers were built in Lean 4.28.0 against Mathlib revision
\texttt{8f9d9cff6bd728b17a24e163c9402775d9e6a365}.  The submitted project reports a successful \texttt{lake build}, no \texttt{sorry}, no \texttt{admit}, and no custom \texttt{axiom} declarations.  The principal checked theorems depend only on the standard Lean/Mathlib logical infrastructure \texttt{propext}, \texttt{Classical.choice}, and \texttt{Quot.sound}.  These facts should not be confused with a claim that the development is constructive in the strict foundational sense.

\subsection{Why the title says ``Dynamics'' although no dynamics is present}

The current paper does \emph{not} derive inertia, an equation of motion, backreaction, or an interaction law.  The use of ``Dynamics'' names the research programme rather than the present result.  The result of Tasks 1 and 2 is that the available kinematic structures split into branches whose mutual relation is not derived.  The missing cross-branch relation is therefore precisely where any later dynamical proposal would have to become mathematically explicit.

\section{Heuristic origin and its mathematical quarantine}
\label{sec:heuristic}

The internal visualization that generated the questions was intentionally elementary.  Imagine two oppositely directed periodic carriers, with marks at regular intervals, together with a local periodic process associated with a timelike object.  One can then ask, purely heuristically, whether an observed change should be pictured as a change of mark spacing, a change of pattern rate, a combination of both, or a change of the local periodic rate.

This picture is \emph{not} used as ontology and is not supplied to the theorem prover.  Its only role is to generate candidate mathematical distinctions.  The translation used to eliminate the picture is summarized in \cref{tab:translation}.

\begin{table}[ht]
\centering
\caption{Internal heuristic vocabulary and the mathematical objects used in the formal audit.  The left column is explanatory only; all proofs use only the right column.}
\label{tab:translation}
\begin{tabular}{>{\raggedright\arraybackslash}p{0.34\linewidth} >{\raggedright\arraybackslash}p{0.58\linewidth}}
\toprule
Internal visualization & Mathematical replacement \\
\midrule
Two counter-propagating ``carriers'' & Two oriented Lorentz-null sectors $\sigma\in\{+1,-1\}$; equivalently the complementary split-complex idempotent sectors $e_\pm$.\\
Spacing of periodic marks & Wavenumber $k_\sigma$ or wavelength $\lambda_\sigma=2\pi/|k_\sigma|$ where $k_\sigma\neq0$.\\
Pattern or ``carrier'' rate & Wave frequency $\omega_\sigma$ and the ratio $\omega_\sigma/k_\sigma$; under standard null dispersion $\omega_\sigma=\sigma c k_\sigma$.\\
Local periodic process & A scalar phase lift $\Theta(\tau)=\Theta_0+\Omega\tau$ along a timelike worldline.\\
Local motion & Rapidity $\eta$, coordinate velocity $v=c\tanh\eta$, and proper time $\tau$.\\
Observed meetings / synchronization of marks & Trajectory-local phase rate $d\phi_\sigma/dt=k_\sigma v-\omega_\sigma$.\\
A limiting ``direction'' at $v\to c$ & Positive null ray or projective null direction, not a unit timelike vector at $v=c$.\\
\bottomrule
\end{tabular}
\end{table}

After this translation, the heuristic is retained only as an aid to visualization.  It is not used to interpret the mathematics.  In particular, a change in $k_\sigma$ is not called a physical spatial deformation, a change in $\omega_\sigma$ is not called a changing medium speed, and a quantity tending to zero is not interpreted as something physically stopping.

A further methodological control is that Aristotle was \emph{blind to the heuristic}.  It received standard mathematical objects, comparison questions, and hard-stop conditions only.  Thus any theorem reported below was not obtained by encoding the visualization as a premise.

\section{Task 1: standard kinematics and identifiability}

\subsection{Split-complex and Lorentz structure}

Let
\begin{equation}
\D=\{a+jb:a,b\in\R,\quad j^2=+1\}
\end{equation}
with complementary idempotents
\begin{equation}
 e_\pm=\frac{1\pm j}{2},\qquad e_\pm^2=e_\pm,\qquad e_+e_-=0.
\end{equation}
The Lean development proves a ring isomorphism
\begin{equation}
 a+jb\longmapsto (a+b,a-b)
\end{equation}
between $\D$ and $\R\times\R$ with componentwise multiplication.  Historically, the algebra belongs to the family introduced by Cockle in his work on tessarines and related hyperbolic units \cite{Cockle1849New,Cockle1849Systems}.  Here no historical construction is imported as a proof premise; the required ring facts are proved in Lean.

On $V=\R^{1,1}$ the convention is
\begin{equation}
 Q_c(t,x)=c^2t^2-x^2,
\end{equation}
with null coordinates
\begin{equation}
 u=ct+x,\qquad w=ct-x,
\end{equation}
so that
\begin{equation}
 Q_c(t,x)=uw.
\end{equation}
The standard proper orthochronous boost is written in rapidity form.  In null coordinates its action is
\begin{equation}
 (u,w)\longmapsto(e^{\eta}u,e^{-\eta}w).
\label{eq:nullboost}
\end{equation}
The formal development derives this reciprocal scaling from form preservation, determinant $1$, and orthochronous orientation.  It also proves that the same factors are the eigenvalues on the two split-complex idempotent sectors.  This is the familiar Lorentz geometry originating in special relativity \cite{Einstein1905,Minkowski1909}, expressed in a split-complex basis.

The rapidity/velocity relation is
\begin{equation}
 \beta=\frac vc=\tanh\eta,
 \qquad
 \gamma=\cosh\eta,
 \qquad
 \frac{d\tau}{dt}=\sqrt{1-\beta^2}=\sech\eta.
\label{eq:properrate}
\end{equation}
Task 1 explicitly separates the reciprocal null-sector factors $(e^{\eta},e^{-\eta})$ from $(\tanh\eta,\sech\eta)$.  They are functions of the same rapidity but are not the same pair; for example,
\begin{equation}
 \tanh\eta=\frac{e^{\eta}-e^{-\eta}}{e^{\eta}+e^{-\eta}},
 \qquad
 \sech\eta=\frac{2}{e^{\eta}+e^{-\eta}}.
\end{equation}

\subsection{Periodic null-sector phases}

For each oriented sector introduce
\begin{equation}
 \phi_\sigma(t,x)=k_\sigma x-\omega_\sigma t.
\label{eq:phase}
\end{equation}
Task 1 does not assume the transformed wave covector.  It derives the transformation law by imposing phase invariance and proves its uniqueness.  In null components the covector factors transform reciprocally to the null coordinates, as required for invariance of \eqref{eq:phase}.

Along a timelike trajectory $x=x(t)$ with coordinate velocity $v$, the observable phase rate is
\begin{equation}
 \frac{d\phi_\sigma}{dt}=k_\sigma v-\omega_\sigma.
\label{eq:obs}
\end{equation}
Define
\begin{equation}
 \mathcal O_v(k,\omega)=kv-\omega.
\end{equation}
For fixed known $v$ and observed value $r$, Task 1 proves
\begin{equation}
 \mathcal O_v^{-1}(r)
 =\{(s,sv-r):s\in\R\},
\label{eq:fibre}
\end{equation}
so the fibre is a one-dimensional affine line.  Therefore $k$ and $\omega$ are \emph{not} separately identifiable from a single trajectory-local phase rate.  A proper-time phase rate merely multiplies \eqref{eq:obs} by the known factor $\gamma$ and adds no independent information.  By contrast, two distinct known velocities determine $(k,\omega)$ uniquely.

This is the first important negative result: trajectory-local kinematics alone does not decide whether a change in the observed phase rate should be parameterized through $k$, through $\omega$, or through both.

\subsection{Models A, B, and C}

Task 1 compares three parameterizations:
\begin{align}
\text{A:}\quad & \omega_\sigma=\sigma ck_\sigma,\qquad k_\sigma=k_\sigma(v),\\
\text{B:}\quad & k_\sigma=k_{0,\sigma},\qquad \omega_\sigma=\omega_\sigma(v),\\
\text{C:}\quad & k_\sigma=k_\sigma(v),\qquad \omega_\sigma=\omega_\sigma(v).
\end{align}
Without an additional dispersion condition, B can reproduce any prescribed trajectory-local rate by choosing
\begin{equation}
 \omega(v)=k_0v-r(v),
\end{equation}
and C contains an entire functional degree of freedom corresponding to the choice of $k(v)$.

The standard null condition
\begin{equation}
 \omega^2=c^2k^2
\label{eq:dispersion}
\end{equation}
changes the conclusion.  For $k\neq0$,
\begin{equation}
 \left|\frac{\omega}{k}\right|=c.
\end{equation}
Thus a genuinely variable propagation speed at fixed nonzero $k$ is incompatible with the same standard Lorentz-null structure.  Models A and B can therefore be trajectory-locally observationally indistinguishable while being structurally inequivalent once null dispersion is imposed.

The original Task-1 prose described the null-constrained Model-C fibre as ``exactly two solutions''.  Task 2 later sharpened this: the two algebraic branches are distinct only for nonzero observed rate $r$.  At $r=0$ they coincide, giving exactly one distinct solution; for $r\neq0$ there are exactly two.

\subsection{Vanishing rates: endpoint behaviour does not identify a quantity}

Task 1 defines
\begin{equation}
 R_\tau(v)=c\frac{d\tau}{dt}=c\sqrt{1-\frac{v^2}{c^2}}=c\sech\eta,
\end{equation}
with
\begin{equation}
 R_\tau(0)=c,
 \qquad
 \lim_{v\to c^-}R_\tau(v)=0,
 \qquad
 v^2+R_\tau(v)^2=c^2.
\end{equation}
It also proves that $R_\tau$ is a bijective coordinate on the speed interval $[0,c)$.  The first written summary overstated what follows from this factorization property.  Task 2 corrects the point: factorization
\begin{equation}
 F(v)=f(R_\tau(v))
\end{equation}
does not imply $F=R_\tau$ and does not canonically select $R_\tau$.

Indeed, several natural standard functions have the same endpoints, including the closing-speed quantity and the relativistic Doppler factor multiplied by $c$.  Task 2 formally proves that the selected examples are pairwise distinct at $v=c/2$.  Hence
\begin{equation}
 R(0)=c,
 \qquad
 R(v)\to0\quad(v\to c^-)
\end{equation}
by itself identifies no unique standard scalar rate.

\subsection{Massive and null structures}

Task 1 keeps the genuine null case separate from the limit of massive motion.  For $m>0$,
\begin{equation}
 u^\mu=c(\cosh\eta,\sinh\eta),
 \qquad
 p^\mu=mc(\cosh\eta,\sinh\eta),
\end{equation}
with
\begin{equation}
 E^2-c^2p^2=m^2c^4.
\end{equation}
At fixed $v<c$, the limit $m\to0^+$ sends $E$ and $p$ to zero and does not yield a nonzero null momentum.  At fixed nonzero energy, $\gamma\to\infty$, $|\eta|\to\infty$, and $v\to c$, while the four-momentum tends to a nonzero null momentum.

For a genuine null worldline, Lorentz arclength remains mathematically defined and is identically zero.  Task 2 corrects the earlier wording ``proper time ceases to exist'': the arclength degenerates and cannot be used as a nondegenerate worldline parameter.  An affine parameter can be used instead, but its scale is not canonically normalized by the null geometry alone.

\section{Task 2: carrier dependence and the projective null boundary}

Task 1 compared timelike and null \emph{vectors}.  That comparison left a subtle ambiguity: a null vector contains a positive scale that a null \emph{direction} does not.  Task 2 therefore distinguishes three carriers for a nonzero $X\in V$:
\begin{equation}
 X,
 \qquad
 \Rayp(X)=\{\lambda X:\lambda>0\},
 \qquad
 [X]\in\mathbb P(V).
\end{equation}
Mathlib's standard projectivization identifies nonzero vectors up to nonzero scalar multiplication \cite{MathlibProjectivization}.  On the future cone, where the time orientation fixes the sign, Task 2 proves that equality of positive rays and equality of projective classes coincide.

\subsection{The projectivized future causal cone}

Every future causal ray has a normalized representative of the form
\begin{equation}
 (1,\beta),
 \qquad -1\le\beta\le1,
\end{equation}
after the fixed $c$-normalization.  Task 2 proves a bijection
\begin{equation}
 \{\text{future causal rays}\}\cong[-1,1],
\end{equation}
with
\begin{equation}
 |\beta|<1
 \quad\Longleftrightarrow\quad
 \text{timelike direction},
\end{equation}
and
\begin{equation}
 \beta=\pm1
 \quad\Longleftrightarrow\quad
 \text{null direction}.
\end{equation}

The induced connected boost action is the fractional-linear map
\begin{equation}
 F_\theta(\beta)
 =\frac{\beta+\tanh\theta}{1+\beta\tanh\theta},
\label{eq:Ftheta}
\end{equation}
and Task 2 proves
\begin{equation}
 F_{\theta_1}\circ F_{\theta_2}=F_{\theta_1+\theta_2}.
\end{equation}
The interior $(-1,1)$ is one transitive orbit, while each endpoint is fixed by every connected boost:
\begin{equation}
 F_\theta(+1)=+1,
 \qquad
 F_\theta(-1)=-1.
\end{equation}
Spatial parity exchanges the two endpoints.

For the timelike unit direction,
\begin{equation}
 U(\eta)=c(\cosh\eta,\sinh\eta),
\end{equation}
the scale-free representative is $(1,\tanh\eta)$.  Hence its direction coordinate tends to $\pm1$ as $\eta\to\pm\infty$.  The four-velocity itself diverges componentwise; it is the projective/ray direction that has a finite boundary limit.  In the formal development this boundary result is established through the canonical $\beta$-coordinate and the proved ray--interval bijection; a fully topologized homeomorphism statement inside Mathlib's projective-space topology would be a possible further hardening, but is not required for the orbit/stabilizer conclusions below.

\subsection{The null stabilizer changes when scale is removed}

The sharpest Task-2 result is carrier dependent.  Let $K_\sigma(E)$ be a nonzero future null vector in sector $\sigma$, with positive scale $E$.  The boost acts as
\begin{equation}
 B_\theta K_\sigma(E)=K_\sigma(e^{\sigma\theta}E).
\label{eq:nullscale}
\end{equation}
Task 2 proves both inclusions in
\begin{equation}
 \Orb_{\mathrm{vec}}(K_\sigma(E))
 =\Rayp(K_\sigma(E)).
\label{eq:orbitray}
\end{equation}
Thus the complete positive null ray is already the orbit of one scaled null vector.

For a \emph{fixed vector}, preserving the scale forces $\theta=0$:
\begin{equation}
 \Stab_{\mathrm{vec}}(K_\sigma(E))=\{0\}.
\end{equation}
After quotienting by positive scale, however, every boost maps the vector to another representative of the same ray.  Consequently
\begin{equation}
 \Stab_{\mathrm{ray}}(\Rayp(K_\sigma(E)))=\R,
 \qquad
 \Stab_{\mathrm{proj}}([K_\sigma(E)])=\R.
\label{eq:stabjump}
\end{equation}
The null orbit therefore collapses to one point in the ray/projective quotient.

The timelike case behaves differently: quotienting a future timelike vector by positive scale does not enlarge its stabilizer.  Its connected-boost stabilizer remains trivial.  The boundary-symmetry statement from Task 1 must therefore be made carrier-specific:
\begin{equation}
\begin{array}{c|c}
\text{carrier/state} & \text{connected-boost stabilizer}\\
\hline
\text{timelike unit vector} & \{0\}\\
\text{timelike positive ray} & \{0\}\\
\text{fixed nonzero null vector} & \{0\}\\
\text{positive null ray} & \R\\
\text{projective null direction} & \R
\end{array}
\end{equation}
This isotropy change is not new physics.  It is a quotient-dependent fact of standard $1+1$ Lorentz geometry: once the scale that the boost changes in \eqref{eq:nullscale} is declared unobservable at the level of the carrier, the boost can no longer move the resulting projective direction.

\section{Task 2: an independent scalar phase on a timelike curve}

The second Task-2 branch is introduced independently of the null-wave sector.  Let
\begin{equation}
 \Theta(\tau)=\Theta_0+\Omega\tau,
\label{eq:scalarphase}
\end{equation}
where $\Theta$ is defined to be a scalar phase lift and $\Omega$ is a constant real parameter.  No equation relating $\Omega$ to $k_\sigma$ or $\omega_\sigma$ is assumed.

With respect to proper time,
\begin{equation}
 \frac{d\Theta}{d\tau}=\Omega.
\label{eq:intrinsicrate}
\end{equation}
Because $\Theta$ was introduced as a scalar phase and $\tau$ as proper time, $\Omega$ is scalar by construction in this branch; this should not be misread as an independently derived new invariance principle.

For constant rapidity, coordinate time satisfies \eqref{eq:properrate}, hence
\begin{equation}
 \frac{d\Theta}{dt}
 =\Omega\frac{d\tau}{dt}
 =\Omega\sech\eta.
\label{eq:coordrate}
\end{equation}
Thus
\begin{equation}
 \left.\frac{d\Theta}{dt}\right|_{\eta=0}=\Omega,
 \qquad
 \lim_{|\eta|\to\infty}\frac{d\Theta}{dt}=0,
\end{equation}
while \eqref{eq:intrinsicrate} remains constant.

This vanishing coordinate-time phase rate must be kept distinct from the null-sector boost weights.  Task 2 proves
\begin{equation}
 \sech\eta=e^{\eta}\iff\eta=0,
 \qquad
 \sech\eta=e^{-\eta}\iff\eta=0.
\end{equation}
The null-wave scales transform with weights $e^{\sigma\theta}$, whereas the coordinate-time scalar-phase rate contains $\sech\eta$.  Common dimensions or endpoint behaviour therefore do not create an identity between the branches.

\section{The dependency split revealed by Tasks 1 and 2}
\label{sec:missing}

The present mathematical structure can be displayed as three branches:
\begin{equation}
\boxed{
\begin{array}{ccc}
& \text{Lorentz causal/projective geometry} & \\
& \eta,\ B_\eta,\ \beta,\ [K_\pm] & \\
\swarrow & & \searrow\\[-0.3em]
\text{timelike scalar phase} & & \text{null wave-covector sector}\\
\Theta(\tau)=\Theta_0+\Omega\tau & & \omega_\sigma=\sigma ck_\sigma\\
\rho_t=\Omega\sech\eta & & k_\sigma\mapsto e^{\sigma\theta}k_\sigma
\end{array}}
\label{eq:branches}
\end{equation}
The vertical/diagonal dependencies are established.  The horizontal dependency is not:
\begin{equation}
 \boxed{\Omega\ \not\xleftrightarrow[\text{derived}]{\text{current inputs}}\ k_\sigma.}
\label{eq:missingedge}
\end{equation}
Equation \eqref{eq:missingedge} is not an assertion that no relation can ever exist.  It states that no such relation follows from the structures currently assumed and proved.

Task 2 supports this conclusion by an explicit continuous rescaling freedom.  In the positive-scale sector, independent rescalings
\begin{equation}
 \Omega\mapsto a\Omega,
 \qquad
 k_\sigma\mapsto b_\sigma k_\sigma,
 \qquad a,b_\sigma>0,
\end{equation}
preserve the separately stated standard relations.  Hence an equality such as $\Omega=ck_\sigma$ is not invariant under the available independent scale freedoms and cannot be canonical without additional input.

The precise scope matters.  The formal theorem establishes independence of the \emph{positive continuous scales}.  It does not by itself exclude every conceivable discrete relation involving signs, orientation, or the sector label $\sigma$.  Such discrete cross-branch structure has not been derived and remains an open mathematical question rather than a negative theorem.

This is the central result of the paper.  The existing Lorentz geometry supplies a common rapidity parameter and known transformation laws, but it does not supply the missing edge.  Any future equation
\begin{equation}
 \mathcal C(\Omega,k_+,k_-,\eta,\ldots)=0
\end{equation}
would therefore have to be classified as one of the following:
\begin{enumerate}[label=(\roman*)]
\item a theorem derived from additional mathematical structure not yet exposed;
\item an independently motivated coupling assumption;
\item a dynamical law;
\item or a redundant reparameterization of existing kinematics.
\end{enumerate}
Simply defining such a relation would not derive it.

\section{What has and has not been established}

\subsection{Established}

The machine-checked results support the following statements.
\begin{enumerate}[label=\arabic*.]
\item The split-complex idempotent basis and null-coordinate basis are equivalent representations of the same $1+1$ Lorentz algebraic structure.
\item A single trajectory-local phase rate determines $kv-\omega$ but not $k$ and $\omega$ separately; its fibres are affine lines.
\item Standard null dispersion $\omega^2=c^2k^2$ removes the continuous phase-speed freedom and enforces $|\omega/k|=c$ for $k\neq0$.
\item Observational equivalence of Models A and B along one prescribed timelike trajectory does not imply equivalence as Lorentz-null structures.
\item The endpoint conditions $R(0)=c$ and $R(v)\to0$ do not uniquely identify $R_\tau=c\,d\tau/dt$ or any other single standard rate.
\item The genuine null case is not obtained by substituting $v=c$ into all massive formulas.  Null Lorentz arclength is defined but degenerate, and affine null parameters have no canonical scale from the null geometry alone.
\item The projectivized future causal directions form an interval with timelike interior and two null boundary directions.
\item A scaled nonzero null vector has trivial connected-boost stabilizer, but its positive ray/projective direction has the entire connected boost group as stabilizer.
\item An independently defined scalar phase satisfies $d\Theta/d\tau=\Omega$ and $d\Theta/dt=\Omega\sech\eta$.
\item No canonical equation relating the positive continuous scalar-phase scale $\Omega$ to the positive continuous null-wave scale follows from the present assumptions.
\end{enumerate}

\subsection{Not established}

The present work does \emph{not} establish:
\begin{enumerate}[label=\arabic*.]
\item that the scalar phase is a physical clock rather than an abstract scalar periodic degree of freedom;
\item that $k_\sigma$ represents a physical spatial lattice or deformation;
\item that $\omega_\sigma/k_\sigma$ represents a variable medium velocity;
\item that the vanishing of $\Omega\sech\eta$ selects a unique physical mechanism;
\item that $\Omega$ is determined by $k_+$ and $k_-$, or conversely;
\item that the null-ray isotropy jump has a dynamical consequence;
\item an inertial law, interaction law, backreaction law, gravitational interpretation, or new field equation.
\end{enumerate}

\section{Interpretive status of the original visualization}

It is useful to return once, and only once, to the visualization of \cref{sec:heuristic} in order to state what the formal results permit us to say about it.

The mathematics does not support treating the two null sectors as carriers whose standard null propagation speed continuously slows below $c$.  Under \eqref{eq:dispersion}, changes of $k_\sigma$ and $\omega_\sigma$ are linked while $|\omega_\sigma/k_\sigma|=c$ remains fixed.  Separately, the coordinate-time rate of the independent scalar phase can tend to zero through the factor $\sech\eta$ while its proper-parameter rate remains constant.  These facts make it possible to visualize two different kinematic mechanisms, but they do not identify either visualization with physical reality.

Accordingly, from this point onward the heuristic picture is retained only for exposition.  It is not permitted to supply premises, select a preferred mathematical carrier, or identify the missing coupling.  Any later coupling must be stated directly in mathematical form and must carry its own dependency justification.

\section{Open structure and next mathematical question}

The unresolved structure is now sharply localized.  There is a Lorentz/null branch with oriented projective boundary directions and null-wave scale parameters.  There is an independent timelike scalar-phase branch with intrinsic parameter $\Omega$ and coordinate-time rate $\Omega\sech\eta$.  The standard Lorentz group acts on both branches, but the present theory contains no derived map
\begin{equation}
 \text{scalar-phase state}\longleftrightarrow\text{null-sector state}.
\end{equation}

The next mathematically legitimate question is therefore not ``what physical mechanism couples them?'' but:
\begin{quote}
What additional mathematical structure, if any, can define a non-arbitrary relation between the timelike scalar-phase degree of freedom and the oriented null-sector degrees of freedom while preserving the already proved Lorentz and projective structures?
\end{quote}
A negative answer remains admissible.  If no canonical relation exists, the branches are simply independent.  If a relation does exist, the dependency audit must identify exactly which new premise makes it possible.  Only after that point would a discussion of persistence, interaction, or dynamics be mathematically justified.

\section{Formal verification status and limitations}

The Aristotle artifacts for Tasks 1 and 2 contain the corresponding Lean sources and audit documents.  Task 2 imports Task 1 without modifying its source.  According to the generated build reports, the complete project builds under Lean 4.28.0 with the Mathlib revision stated above and contains no placeholders or custom axioms.

Several qualifications are important.
\begin{itemize}
\item The projective boundary theorem is formalized via a canonical ray coordinate $\beta\in[-1,1]$, together with equality results connecting future positive rays and projective classes.  A stronger topological packaging as a Mathlib \texttt{Homeomorph} of the relevant projective subspace is not part of the present task.
\item The continuous-scale independence theorem concerns the positive scale sector.  It should not be cited as excluding every possible discrete relation involving orientation or sector sign.
\item $\Omega$ is invariant in the scalar-phase branch because the phase was introduced as a scalar function of proper time.  This is a carrier choice, not an independently derived physical principle.
\item The term ``dynamics'' in the title does not imply that a dynamical equation has been obtained.  The principal result is the exact location at which such additional structure would have to enter.
\end{itemize}

\section{Conclusion}

Tasks 1 and 2 convert an initially ambiguous visualization into a sharply separated mathematical problem.  Task 1 shows that trajectory-local phase observations are non-identifying before a dispersion relation is supplied, while standard null dispersion restores the fixed null propagation speed and distinguishes observational parameterization from Lorentz-null structure.  Task 2 shows that the null boundary depends essentially on the carrier: the scale-bearing null vector has a nontrivial boost orbit with trivial stabilizer, whereas its scale-free ray/projective class is a boost fixed point with maximal connected-boost stabilizer.  Independently, a scalar phase along a timelike curve possesses a constant proper-parameter rate and a coordinate-time rate proportional to $\sech\eta$.

The resulting dependency graph contains two mathematically controlled kinematic branches but no derived coupling between their continuous scales.  That absence is not a defect to be hidden by definition; it is the present research result.  The first genuinely new structure required by the programme called \emph{Split-Complex Lorentzian Dynamics} is therefore a justified cross-branch relation---or a proof that no canonical such relation exists.

\appendix
\section{Formal theorem traceability}
\label{app:trace}

The following table records the principal Lean declarations supporting the claims emphasized in the paper.  It is not a complete theorem inventory; the complete dependency and obligation ledgers are contained in the Aristotle Task-1 and Task-2 audit artifacts.

\begin{longtable}{>{\raggedright\arraybackslash}p{0.31\linewidth} >{\raggedright\arraybackslash}p{0.61\linewidth}}
\caption{Selected theorem-to-claim traceability.}\label{tab:leantrace}\\
\toprule
Mathematical claim & Lean declaration(s) \\
\midrule
\endfirsthead
\toprule
Mathematical claim & Lean declaration(s) \\
\midrule
\endhead
Split-complex/null-coordinate ring equivalence & \texttt{Task01.SC.toNull}; Lorentz/split norm correspondence \texttt{N\_toSC}.\\
Reciprocal null-coordinate boost scaling & \texttt{nullU\_boost}, \texttt{nullW\_boost}; sector eigenvalues \texttt{sexp\_mul\_ePlus}, \texttt{sexp\_mul\_eMinus}.\\
Rapidity/velocity correspondence & \texttt{vel\_rap}, \texttt{rap\_vel}; proper-time rate \texttt{properTimeRate\_vel}.\\
Wave-covector law from phase invariance & \texttt{phase\_boost\_invariant}, \texttt{waveCovector\_unique}.\\
Single-trajectory non-identifiability & \texttt{obs\_fibre}, \texttt{obs\_not\_injective}, \texttt{obs\_fibre\_eq\_line}; recovery from two velocities \texttt{two\_velocities\_identify}.\\
Null dispersion and fixed null speed & \texttt{nullDispersion\_iff}, \texttt{null\_phase\_speed}, \texttt{null\_group\_velocity}; Model-B rigidity \texttt{modelB\_null\_const}, \texttt{modelB\_no\_variable\_speed}.\\
Task-1 Model-C cardinality repair & \texttt{Task02.modelCNullSolutions\_zero\_ncard}, \texttt{Task02.modelCNullSolutions\_ncard}.\\
Endpoint data do not select a unique rate & \texttt{Task02.endpoint\_data\_do\_not\_select}, \texttt{Task02.vanishing\_rates\_pairwise\_distinct}.\\
Null arclength is defined but degenerate & \texttt{Task02.null\_arclength\_defined\_and\_zero}, \texttt{Task02.null\_arclength\_not\_injective}, \texttt{Task02.null\_affine\_no\_canonical\_scale}.\\
Future causal rays parameterized by $[-1,1]$ & \texttt{Task02.bijOn\_causalRays}, with timelike/null characterizations \texttt{mem\_FutureTimelike\_iff}, \texttt{mem\_FutureNull\_iff}.\\
Fractional-linear projective boost action & \texttt{Task02.betaOf\_boostV}, group law \texttt{Task02.Fmap\_Fmap}; endpoint fixed points \texttt{Fmap\_one}, \texttt{Fmap\_neg\_one}.\\
Null-vector orbit equals its positive ray & \texttt{Task02.orbVec\_nullMomentum}; quotient orbit collapse \texttt{Task02.null\_orbit\_collapses}.\\
Carrier-dependent stabilizer jump & \texttt{Task02.stabVec\_nullMomentum}, \texttt{Task02.stabRay\_nullMomentum}, \texttt{Task02.stabProj\_nullMomentum}; consolidated statement \texttt{Task02.boundary\_symmetry\_carrier\_specific}.\\
Scalar-phase rates & \texttt{Task02.hasDerivAt\_scalarPhase\_properTime}, \texttt{Task02.hasDerivAt\_scalarPhase\_coordTime}.\\
Continuous positive-scale cross-branch independence & \texttt{Task02.no\_canonical\_cross\_sector\_relation}, \texttt{Task02.cross\_sector\_equality\_not\_invariant}.\\
$\sech\eta$ is not a null-sector exponential weight except at rest & \texttt{Task02.sech\_eq\_exp\_iff}, \texttt{Task02.sech\_eq\_exp\_neg\_iff}.\\
\bottomrule
\end{longtable}

\section{Reproducibility metadata}

The Task-2 artifact records the following verification metadata:
\begin{itemize}
\item Lean toolchain: \texttt{leanprover/lean4:v4.28.0};
\item Mathlib revision: \texttt{8f9d9cff6bd728b17a24e163c9402775d9e6a365};
\item Task-1 modules: \texttt{SplitComplex.lean}, \texttt{Boost.lean}, \texttt{ProperTime.lean}, \texttt{Phase.lean}, \texttt{Boundary.lean};
\item Task-2 modules: \texttt{Carriers.lean}, \texttt{ProjectiveCone.lean}, \texttt{Orbits.lean}, \texttt{ScalarPhase.lean}, \texttt{Task01Audit.lean};
\item reported full build: successful;
\item placeholders: zero \texttt{sorry}, zero \texttt{admit}, zero custom \texttt{axiom} declarations;
\item principal Task-2 theorem axiom reports: only \texttt{propext}, \texttt{Classical.choice}, and \texttt{Quot.sound}.
\end{itemize}

\section*{Acknowledgements and provenance}

The mathematical prompts were prepared so that Aristotle received no access to the internal visualization described in \cref{sec:heuristic}.  The purpose of this separation was to reduce confirmation bias: heuristic language generated questions but was not allowed to become a theorem premise.  Aristotle produced the Lean formalizations and machine-generated audit reports for Tasks 1 and 2.  ChatGPT assisted with critical comparison of the formal statements, identification of over-broad prose claims in the first audit, and preparation of this manuscript.  Responsibility for the research framing and final editorial decisions remains with the project editor.

\begin{thebibliography}{99}

\bibitem{Cockle1849New}
J. Cockle,
``On a New Imaginary in Algebra,''
\emph{The London, Edinburgh, and Dublin Philosophical Magazine and Journal of Science}, Series 3, vol. 34, pp. 37--47, 1849.
DOI: \href{https://doi.org/10.1080/14786444908646169}{10.1080/14786444908646169}.

\bibitem{Cockle1849Systems}
J. Cockle,
``On Systems of Algebra Involving More than One Imaginary; and on Equations of the Fifth Degree,''
\emph{The London, Edinburgh, and Dublin Philosophical Magazine and Journal of Science}, Series 3, vol. 35, no. 238, pp. 434--437, 1849.
DOI: \href{https://doi.org/10.1080/14786444908646384}{10.1080/14786444908646384}.

\bibitem{Einstein1905}
A. Einstein,
``Zur Elektrodynamik bewegter K\"orper,''
\emph{Annalen der Physik}, vol. 322, no. 10, pp. 891--921, 1905.
DOI: \href{https://doi.org/10.1002/andp.19053221004}{10.1002/andp.19053221004}.

\bibitem{Minkowski1909}
H. Minkowski,
``Raum und Zeit,''
address delivered in Cologne, 21 September 1908; published in
\emph{Physikalische Zeitschrift}, vol. 10, pp. 104--111, 1909.

\bibitem{ONeill1983}
B. O'Neill,
\emph{Semi-Riemannian Geometry with Applications to Relativity},
Academic Press, 1983.

\bibitem{PenroseRindler1984}
R. Penrose and W. Rindler,
\emph{Spinors and Space-Time, Volume 1: Two-Spinor Calculus and Relativistic Fields},
Cambridge University Press, 1984.

\bibitem{MathlibProjectivization}
The Mathlib Community,
``Mathlib.LinearAlgebra.Projectivization.Basic,''
official Mathlib 4 documentation and source API.
\url{https://leanprover-community.github.io/mathlib4_docs/Mathlib/LinearAlgebra/Projectivization/Basic.html}.
The formal artifact used Mathlib revision \texttt{8f9d9cff6bd728b17a24e163c9402775d9e6a365}.

\bibitem{AristotleArtifact}
Aristotle / Harmonic,
``Task 01--02 formalization request: Foundation, Identifiability, Projective Causal Boundary, and Scalar-Phase Audit,'' 2026.
\url{https://aristotle.harmonic.fun/dashboard/requests/201d2dcd-4452-4ae1-8656-4637890c2223}.

\bibitem{PreviousExperiment}
J. Seeck et al.,
``Split-Complex Lorentzian Kinematics,'' project discussion and previous experiment.
\url{https://www.astronews.com/community/threads/split-complex-lorentzian-kinematics.12707/post-156281}.

\end{thebibliography}

\end{document}
